%======================================================================
\section{Energy levels and the stability of atoms}\label{sec:levels}
%======================================================================

\begin{objectives}
  You will see why the energy of a confined particle comes in discrete levels, why a system in a single level does not change while a superposition of levels oscillates at the Bohr frequencies, how this explains sharp spectral lines, and why the uncertainty principle stops the electron of a hydrogen atom from falling into the nucleus.
\end{objectives}

\begin{recall}
  The spectral theorem and the functional calculus (Section~\ref{sec:operators}); the Schr\"odinger equation and energy eigenstates (Proposition~\ref{prop:schrodinger}); the uncertainty principle (Theorem~\ref{thm:uncertainty}) and the Gaussian wave packet (Example~\ref{ex:gaussian-uncertainty}); wave functions and the Born rule (Example~\ref{ex:box}).
\end{recall}

The second and third puzzles of the Introduction are about energy. Hydrogen emits light only at a few sharp frequencies (Figure~\ref{fig:levels}), and the atom does not collapse although classical physics says it must. Both are explained by one fact of the theory: the possible energies of a system are the eigenvalues of its Hamiltonian.

\subsection{Stationary states and Bohr frequencies}

Let $H=\sum_nE_n\ketbra{E_n}{E_n}$ be the spectral decomposition of the Hamiltonian of a system (Theorem~\ref{thm:spectral}), with $E_1<E_2<\cdots$ and, for simplicity, one eigenvector per eigenvalue. We saw after Proposition~\ref{prop:schrodinger} that a system in one energy eigenstate only picks up a global phase, $e^{-\ii Ht/\hbar}\ket{E_n}=e^{-\ii E_nt/\hbar}\ket{E_n}$, so nothing observable about it ever changes: it is a \term{stationary state}. A superposition of levels behaves differently.

\begin{proposition}[Only Bohr frequencies appear]\label{prop:bohr}
  Let $\ket{\psi(0)}=\sum_nc_n\ket{E_n}$ and let $A$ be any observable, with matrix elements $A_{mn}=\bra{E_m}A\ket{E_n}$. Then
  \[
    \langle A\rangle_{\psi(t)}=\sum_{m,n}\overline{c_m}\,c_n\,A_{mn}\,e^{\ii(E_m-E_n)t/\hbar}.
  \]
  The average of every observable is a combination of oscillations at the angular frequencies $\omega_{mn}=(E_m-E_n)/\hbar$, the \term{Bohr frequencies}, and of no others.
\end{proposition}
\begin{proof}
  \begin{steps}
    \item \emph{Evolve} (functional calculus): $\ket{\psi(t)}=e^{-\ii Ht/\hbar}\ket{\psi(0)}=\sum_nc_ne^{-\ii E_nt/\hbar}\ket{E_n}$.
    \item \emph{The bra} is $\bra{\psi(t)}=\sum_m\overline{c_m}e^{\ii E_mt/\hbar}\bra{E_m}$ (conjugate-linearity).
    \item \emph{Sandwich $A$:} $\bra{\psi(t)}A\ket{\psi(t)}=\sum_{m,n}\overline{c_m}c_n\,e^{\ii(E_m-E_n)t/\hbar}\bra{E_m}A\ket{E_n}$. \qedhere
  \end{steps}
\end{proof}

The diagonal terms $m=n$ are constant; only the cross terms move, and each one moves at the frequency of a \emph{difference} of two levels. This is the relative phase at work again: between two levels it grows like $(E_m-E_n)t/\hbar$, just as the relative phase between the two paths of an interferometer grew with $\varphi$.

\paragraph{Spectral lines.}
An electron is a charge. If its state is a superposition of two levels, its average position oscillates at a Bohr frequency, and an oscillating charge radiates light of that frequency, just as an antenna does. Hence the atom emits, and absorbs, light only at the frequencies
\[
  \nu_{mn}=\frac{\omega_{mn}}{2\pi}=\frac{E_m-E_n}{h},\qquad\text{that is,}\qquad h\nu_{mn}=E_m-E_n .
\]
This is Bohr's rule: the energy $h\nu$ of the light emitted equals the energy lost by the atom. A discrete set of levels gives a discrete set of lines. (A full theory of emission treats the light as a quantum system too, and Section~\ref{sec:blackbody} takes a first step in that direction; the frequencies, however, are already fixed by Proposition~\ref{prop:bohr}.)

\begin{example}[The red line of hydrogen]\label{ex:hydrogen-line}
  The levels of hydrogen are $E_n=-13.6\,\mathrm{eV}/n^2$. Find the energy and the wavelength of the light emitted when the electron drops from $n=3$ to $n=2$, using $E=hc/\lambda$ with $hc\approx1240\ \mathrm{eV\,nm}$.
\end{example}
\begin{solution}
  \begin{steps}
    \item \emph{The levels.} $E_3=-13.6/9\approx-1.51$ eV and $E_2=-13.6/4=-3.40$ eV.
    \item \emph{The energy of the light.} $h\nu=E_3-E_2\approx1.89$ eV.
    \item \emph{The wavelength.} $\lambda=hc/(h\nu)\approx1240/1.89\approx656$ nm.
  \end{steps}
  This is the bright red line of hydrogen, and the drops from $n=4,5,\dots$ to $n=2$ give the other visible lines. One formula for the levels explains the whole spectrum; what remains is to understand why the levels are discrete at all.
\end{solution}

\subsection{The particle in a box}

The simplest system with discrete levels is a particle of mass $m$ that moves freely on a segment $[0,L]$ and cannot leave it (a ``box''). Its state is a wave function $\psi(x)$ with $\psi(0)=\psi(L)=0$, and its Hamiltonian is the kinetic energy
\[
  H=\frac{\hat p^2}{2m}=-\frac{\hbar^2}{2m}\frac{d^2}{dx^2},\qquad\hat p=-\ii\hbar\frac{d}{dx}
\]
(Exercise~\ref{ex:no-xp} introduced $\hat p$). The space of wave functions is infinite-dimensional, which Chapter~\ref{ch:tools} did not cover; but we shall see that $H$ has an orthonormal basis of eigenvectors, exactly as the spectral theorem would lead us to expect, and that is all we need.

\begin{theorem}[Levels of a particle in a box]\label{thm:box}
  The eigenvalue problem $H\varphi=E\varphi$ with $\varphi(0)=\varphi(L)=0$ and $\varphi\neq0$ has the solutions
  \[
    E_n=\frac{n^2\pi^2\hbar^2}{2mL^2},\qquad\varphi_n(x)=\sqrt{\frac2L}\,\sin\frac{n\pi x}{L},\qquad n=1,2,3,\dots,
  \]
  and the functions $\varphi_n$ are orthonormal: $\int_0^L\varphi_m\varphi_n\dd x=\delta_{mn}$.
\end{theorem}
\begin{proof}
  \begin{steps}
    \item \emph{Rewrite the equation.} $H\varphi=E\varphi$ means $\varphi''=-k^2\varphi$ with $k^2=2mE/\hbar^2$.
    \item \emph{The energy is positive.} Multiply $-\frac{\hbar^2}{2m}\varphi''=E\varphi$ by $\overline\varphi$ and integrate by parts, using $\varphi(0)=\varphi(L)=0$: $E\int_0^L\abs\varphi^2\dd x=\frac{\hbar^2}{2m}\int_0^L\abs{\varphi'}^2\dd x$. The right side is $>0$ (if $\varphi'=0$ then $\varphi$ is constant, hence $0$), so $E>0$ and $k>0$.
    \item \emph{General solution.} The solutions of $\varphi''=-k^2\varphi$ are $\varphi(x)=a\sin kx+b\cos kx$.
    \item \emph{Left wall.} $\varphi(0)=b=0$.
    \item \emph{Right wall.} $\varphi(L)=a\sin kL=0$ with $a\neq0$ forces $kL=n\pi$ for an integer $n\ge1$. Then $E=\hbar^2k^2/2m=n^2\pi^2\hbar^2/2mL^2$.
    \item \emph{Normalise.} $\int_0^L\sin^2\frac{n\pi x}{L}\dd x=\frac L2$, so $a=\sqrt{2/L}$.
    \item \emph{Orthogonality.} For $m\neq n$, $2\sin\alpha\sin\beta=\cos(\alpha-\beta)-\cos(\alpha+\beta)$ gives $\int_0^L\varphi_m\varphi_n\dd x=\frac1L\int_0^L\bigl(\cos\frac{(m-n)\pi x}L-\cos\frac{(m+n)\pi x}L\bigr)\dd x=0$. \qedhere
  \end{steps}
\end{proof}

\begin{figure}[htb]
  \centering
  \begin{tikzpicture}
    \begin{axis}[width=0.62\linewidth, height=6.2cm, xmin=-0.08, xmax=1.08, ymin=-0.3, ymax=10.6,
        axis lines=none, clip=false]
      \draw[ink, very thick] (axis cs:0,-0.2) -- (axis cs:0,10.4);
      \draw[ink, very thick] (axis cs:1,-0.2) -- (axis cs:1,10.4);
      \draw[ink!60, dashed] (axis cs:0,1) -- (axis cs:1,1);
      \draw[ink!60, dashed] (axis cs:0,4) -- (axis cs:1,4);
      \draw[ink!60, dashed] (axis cs:0,9) -- (axis cs:1,9);
      \node[anchor=west,font=\footnotesize] at (axis cs:1.03,1) {$E_1$};
      \node[anchor=west,font=\footnotesize] at (axis cs:1.03,4) {$E_2=4E_1$};
      \node[anchor=west,font=\footnotesize] at (axis cs:1.03,9) {$E_3=9E_1$};
      \addplot[navy, thick, samples=80, domain=0:1] {1+0.9*sin(deg(pi*x))};
      \addplot[navy, thick, samples=80, domain=0:1] {4+0.9*sin(deg(2*pi*x))};
      \addplot[navy, thick, samples=80, domain=0:1] {9+0.9*sin(deg(3*pi*x))};
      \node[font=\footnotesize,anchor=north] at (axis cs:0,-0.25) {$0$};
      \node[font=\footnotesize,anchor=north] at (axis cs:1,-0.25) {$L$};
    \end{axis}
  \end{tikzpicture}
  \caption{The three lowest levels of a particle in a box. Each eigenfunction $\varphi_n$ is drawn at the height of its energy: it fits $n$ half-wavelengths between the walls, and the energies grow like $n^2$.}
  \label{fig:box}
\end{figure}

The discreteness comes from the boundary conditions. A wave that must vanish at both walls has to fit a whole number of half-wavelengths into the box, as the vibrations of a guitar string do (Figure~\ref{fig:box}); each allowed wavelength $2L/n$ gives one allowed momentum $\hbar k=n\pi\hbar/L$ and one allowed energy. A free particle, with no walls, has no such restriction, and its energies form a continuum (Example~\ref{ex:gaussian-uncertainty} used such free wave functions).

Two further features deserve comment.
\begin{steps}
  \item \emph{The lowest energy is not zero.} $E_1=\pi^2\hbar^2/2mL^2>0$: a confined particle can never be at rest. This \term{zero-point energy} is forced by the uncertainty principle. Confinement to $[0,L]$ makes $\Delta x$ at most of order $L$, so $\Delta p\ge\hbar/2\Delta x$ is at least of order $\hbar/L$, and the average kinetic energy $\langle\hat p^2\rangle/2m\ge(\Delta p)^2/2m$ is at least of order $\hbar^2/mL^2$.
  \item \emph{Smaller boxes have more widely spaced levels.} The levels scale like $1/L^2$.
\end{steps}

\begin{example}[An electron in a nanometre box]\label{ex:nano-box}
  An electron ($m\approx9.11\times10^{-31}$ kg) is confined to a box of length $L=1$ nm. Find $E_1$, the energy of the transition $2\to1$, and the wavelength of the corresponding light. Use $\hbar\approx1.055\times10^{-34}$ J\,s and $1\ \mathrm{eV}\approx1.602\times10^{-19}$ J.
\end{example}
\begin{solution}
  \begin{steps}
    \item \emph{Ground level.} $E_1=\dfrac{\pi^2\hbar^2}{2mL^2}=\dfrac{(3.31\times10^{-34})^2}{2\cdot9.11\times10^{-31}\cdot10^{-18}}\approx6.0\times10^{-20}\ \mathrm J\approx0.38\ \mathrm{eV}$.
    \item \emph{Transition.} $E_2-E_1=(4-1)E_1=3E_1\approx1.13$ eV.
    \item \emph{Wavelength.} $\lambda\approx1240/1.13\approx1100$ nm, in the near infrared.
  \end{steps}
  Energies of an electronvolt for sizes of a nanometre: this is the scale of atoms and of chemistry. Nanocrystals called \emph{quantum dots} are boxes of this kind, and their colour is set by their size: smaller dots have more widely spaced levels and glow bluer.
\end{solution}

\begin{example}[A superposition in the box moves]\label{ex:box-moves}
  Let $\ket{\psi(0)}=\tfrac1{\sqrt2}(\varphi_1+\varphi_2)$. Show that the average position oscillates about the centre of the box at the Bohr frequency $\omega_{21}=(E_2-E_1)/\hbar$, and find the amplitude of the oscillation, given that $\int_0^Lx\,\varphi_1\varphi_2\dd x=-\tfrac{16L}{9\pi^2}$.
\end{example}
\begin{solution}
  \begin{steps}
    \item \emph{Diagonal terms.} $\int_0^Lx\,\varphi_n^2\dd x=\tfrac L2$ for every $n$, because $\varphi_n^2$ is symmetric about $x=\tfrac L2$.
    \item \emph{Apply Proposition~\ref{prop:bohr}} with $A=\hat x$, $c_1=c_2=\tfrac1{\sqrt2}$, and $x_{12}=x_{21}=-\tfrac{16L}{9\pi^2}$ (the eigenfunctions are real):
      \[
        \langle\hat x\rangle(t)=\tfrac12\cdot\tfrac L2+\tfrac12\cdot\tfrac L2+\tfrac12x_{12}\bigl(e^{-\ii\omega_{21}t}+e^{\ii\omega_{21}t}\bigr)=\tfrac L2-\tfrac{16L}{9\pi^2}\cos\omega_{21}t .
      \]
    \item \emph{Amplitude.} $\tfrac{16}{9\pi^2}\approx0.18$: the particle sloshes from side to side with an amplitude of about $18\%$ of the box.
  \end{steps}
  A charged particle doing this radiates at frequency $\omega_{21}/2\pi$, which is the mechanism behind the spectral lines described above. Each stationary state alone gives a constant $\langle\hat x\rangle=L/2$.
\end{solution}

\subsection{Why atoms do not collapse}

The hydrogen atom has the Hamiltonian
\[
  H=\frac{\hat p^2}{2m}-\frac{ke^2}{r},\qquad k=\frac1{4\pi\varepsilon_0},
\]
kinetic energy plus the Coulomb attraction to the nucleus at distance $r$. Classically the electron can lower its energy without limit by falling inwards, since $-ke^2/r\to-\infty$. Quantum mechanically, falling inwards means confining the electron to a smaller region, and confinement costs kinetic energy, as we have just seen. The two effects balance at a definite size.

\begin{example}[Size and energy of the hydrogen atom]\label{ex:hydrogen-estimate}
  Suppose that the electron is spread over a region of radius $r$ around the nucleus. Estimate its energy as a function of $r$, find the radius that minimises it, and evaluate the radius and the energy numerically.
\end{example}
\begin{solution}
  \begin{steps}
    \item \emph{Kinetic energy.} Confinement to size $r$ gives $\Delta x\sim r$, hence $\Delta p\sim\hbar/r$ by the uncertainty principle, and a kinetic energy of order $\hbar^2/2mr^2$ (we drop factors of order one).
    \item \emph{Potential energy.} Of order $-ke^2/r$.
    \item \emph{Total.} $E(r)\approx\dfrac{\hbar^2}{2mr^2}-\dfrac{ke^2}{r}$. The first term wins at small $r$ and the second at large $r$, so $E(r)$ has a minimum.
    \item \emph{Minimise.} $E'(r)=-\dfrac{\hbar^2}{mr^3}+\dfrac{ke^2}{r^2}=0$ gives
      \[
        r=a_0=\frac{\hbar^2}{mke^2},\qquad E(a_0)=\frac{ke^2}{2a_0}-\frac{ke^2}{a_0}=-\frac{ke^2}{2a_0}.
      \]
    \item \emph{Numbers.} With $ke^2\approx2.31\times10^{-28}$ J\,m, $a_0\approx5.3\times10^{-11}$ m, about half an \aa ngstr\"om, and $E(a_0)\approx-2.18\times10^{-18}$ J $\approx-13.6$ eV.
  \end{steps}
  The estimate produces the size of the atom and, by a lucky choice of the factors we dropped, exactly the ground level $E_1=-13.6$ eV of Figure~\ref{fig:levels}. The exact solution of the eigenvalue problem for $H$ gives all the levels $E_n=-13.6\,\mathrm{eV}/n^2$; it is a longer computation, done in every textbook of quantum mechanics~\cite{griffiths2018}.
\end{solution}

So the atom is stable because its Hamiltonian is bounded below: there is a lowest eigenvalue, the ground level, and an atom in its ground state is in a stationary state that has nowhere lower to go. The uncertainty principle, which in Chapter~\ref{ch:postulates} looked like a limitation on what we can know, is here what keeps matter from collapsing.

\begin{lookahead}
  A qubit is, physically, two levels of some system, chosen so that they are well separated from all the others: two levels of an ion, of an atom, or of a superconducting circuit that behaves like an artificial atom. The Bohr frequency between them is the frequency of the light or microwaves used to control them, as the next section shows.
\end{lookahead}

\begin{exerciseset}
  \exercise{1} Find the energy and the wavelength of the light emitted in the transition $n=4\to n=2$ of hydrogen. Is it visible (roughly $400$--$700$ nm)?

  \exercise{1} Verify directly that $\varphi_2(x)=\sqrt{2/L}\,\sin(2\pi x/L)$ satisfies $H\varphi_2=E_2\varphi_2$ and the boundary conditions, and that $\int_0^L\varphi_1\varphi_2\dd x=0$.

  \exercise{2} For the ground state $\varphi_1$ of the box, compute $\langle\hat x\rangle$, $\Delta x$, $\langle\hat p\rangle$ and $\Delta p$. Show that $\Delta x\,\Delta p=\hbar\sqrt{\pi^2/12-1/2}\approx0.57\hbar$, and compare with Theorem~\ref{thm:uncertainty}.

  \exercise{2} Compute $\int_0^Lx\,\varphi_1\varphi_2\dd x=-\tfrac{16L}{9\pi^2}$, the value used in Example~\ref{ex:box-moves}.

  \exercise{2} Estimate the length of a box in which an electron emits red light ($650$ nm) in the transition $2\to1$.

  \exercise{2} Repeat the estimate of Example~\ref{ex:hydrogen-estimate} for an electron bound to a nucleus of charge $Ze$. Show that the radius scales like $1/Z$ and the energy like $Z^2$, and estimate the ground energy of the ion He$^+$ ($Z=2$).

  \exercise{3} Let the spectrum of $H$ be non-degenerate and $\ket{\psi(0)}=\sum_nc_n\ket{E_n}$. Show that the time average $\lim_{T\to\infty}\frac1T\int_0^T\langle A\rangle_{\psi(t)}\dd t$ equals $\sum_n\abs{c_n}^2A_{nn}$, so that over long times only the populations of the levels matter. (Assume the sum is finite.)
\end{exerciseset}
